\chapter{Introdução}

Lesões de pele são áreas anormais ou danificadas da pele, causadas por diversos fatores, como condições ambientais,
infecções ou doenças graves, incluindo o câncer de pele. A identificação precisa da causa é fundamental para o
tratamento adequado, porém, devido à semelhança entre os sintomas, o diagnóstico correto das doenças de pele pode ser
uma tarefa complexa. Além disso, o prognóstico pode se agravar significativamente se o tratamento não for iniciado
cedo o suficiente \cite{habif2015clinical}.

A importância do cuidado com lesões de pele é evidenciada pelo impacto do câncer de pele. No Brasil, 30\% dos tumores
malignos registrados são atribuídos a essa doença \cite{skin_cancer_in_brazil}. Os tipos mais comuns, o carcinoma
basocelular e o carcinoma epidermoide, apresentam baixa letalidade, mas podem deixar sequelas significativas devido ao
tratamento. Já o melanoma, embora menos frequente, é altamente letal e corresponde a cerca de 75\% das mortes causadas
pelo câncer de pele \cite{skin_cancer_screening}. A detecção precoce é essencial para a eficácia do tratamento, pois a
taxa de sobrevida dos pacientes tende a cair com o avanço da doença \cite{skin_cancer_survival}.

A dificuldade da identificação precoce de lesões de pele tem um componente socioeconômico. Em um estudo feito por
\textcite{skin_cancer_socioeconomic}, observa-se que indivíduos com um grau menor de escolaridade têm um prognóstico
pior, pois apresentam tumores em estágios mais avançados. Esta correlação pode indicar que a desigualdade resulta no
atraso do diagnóstico e consequentemente na redução da taxa de sucesso do tratamento.
\textcite{santos2020desigualdades} também constata que a saúde dos brasileiros está associada à renda e classe
econômica.

No Brasil, uma parte da atenção primária é feita por \ac{ACS}, que possuem um treinamento básico para atender a
população e fornecer orientações básicas de saúde. Esses são os profissionais que normalmente estão em contato com as
regiões mais carentes. A identificação da causa de uma lesão de pele está além das responsabilidades de um \ac{ACS}
\cite{filgueiras2011agente}. Isso pode ser problemático, pois a triagem de um caso de lesão de pele requer a análise de
um dermatologista, que nem sempre está disponível em regiões carentes.

Considerando este cenário, seria útil um serviço com boa usabilidade que analisasse imagens de pacientes e recomendasse
uma forma de triagem, de modo a mitigar a carência por dermatologistas. As capacidades de classificar e descrever
imagens com precisão tornam os \acp{MLLM} bons candidatos para fundamentar tal serviço \cite{mllm_success_rate}.

Neste trabalho, o \ac{LLaMA}-3.2-11B foi escolhido como modelo base para o desenvolvimento por ser de código aberto e
por sua relevância no estado da arte para tarefas visuais e textuais \cite{dubey2024llama}. O modelo foi submetido ao
\textit{fine-tuning} com diferentes técnicas de treinamento ao longo de duas etapas, nas quais se utilizaram os
conjuntos de dados \ac{HAM10000} e \ac{STT/SC} para adaptá-lo às tarefas de classificação de lesões de pele e geração
de laudos.

Com o objetivo de utilizar eficientemente os recursos computacionais disponíveis, o \textit{fine-tuning} foi realizado
com técnicas baseadas em \ac{PEFT}. Este método permite a redução do número de parâmetros treinados, resultando em um
menor consumo de recursos \cite{peft}. Especificamente, utilizaram-se os métodos \ac{QLoRA} e \ac{LoRA} para o
treinamento do modelo.

\section{Objetivos}

O objetivo principal deste trabalho é treinar um \ac{LLM} multimodal com métodos de \textit{fine-tuning} para a
classificação de lesões de pele e geração de laudos, avaliando o desempenho dos modelos resultantes e comparando as
diferentes técnicas de treinamento.

\subsection*{Objetivos Específicos}

\begin{itemize}
    \item Realizar o \textit{fine-tuning} do modelo \ac{LLaMA}-3.2-11B com as técnicas de treinamento \ac{QLoRA} e
          \ac{LoRA};
    \item Comparar os modelos resultantes do processo de \textit{fine-tuning} com o \ac{LLaMA}-3.2-11B;
    \item Comparar os modelos resultantes entre si a fim de avaliar o impacto das diferentes técnicas de treinamento.
\end{itemize}

\section{Organização do Trabalho}

O trabalho apresentá inicialmente os conceitos teóricos sobre lesões de pele e \acp{MLLM}, abordando a arquitetura e o
treinamento destes modelos. Em seguida, serão apresentados e discutidos os trabalhos correlatos. Posteriormente, as
duas etapas do desenvolvimento serão detalhadas, incluindo seus aspectos técnicos, treinamentos e resultados. Em
sequência, os dados obtidos serão analisados e avaliados. Por fim, serão apresentadas as conclusões e os possíveis
trabalhos futuros.
